% Generated draft: requires mathematical and source review.
\documentclass[11pt]{article}
\usepackage[margin=25mm]{geometry}
\usepackage{fontspec}
\setmainfont{Latin Modern Roman}
\usepackage{amsmath,amssymb,mathtools,unicode-math}
\setmathfont{Latin Modern Math}
\usepackage{enumitem,needspace,xurl,hyperref}
\setlength{\parindent}{0pt}
\setlength{\parskip}{5pt}
\setlength{\emergencystretch}{3em}
\newcommand{\N}{\mathbb{N}}
\newcommand{\Z}{\mathbb{Z}}
\newcommand{\Q}{\mathbb{Q}}
\newcommand{\R}{\mathbb{R}}
\newcommand{\C}{\mathbb{C}}
\newcommand{\F}{\mathbb{F}}
\newcommand{\A}{\mathbb{A}}
\newcommand{\PP}{\mathbb{P}}
\newcommand{\E}{\mathbb{E}}
\newcommand{\eps}{\varepsilon}
\newcommand{\dd}{\,\mathrm{d}}
\begin{document}
\section*{UnsolvedMath ID: 30002893}
\textbf{Source problem code:} catalogue entry
\subsection{Definitions and mathematical statement}
The supplied record poses a research question about the regularity of polynomial distributions, rather than a fully specified conjecture. Its stated question is: what regularity properties does the distribution \(\mu_Y\) of \(Y=P(X_1,\ldots,X_n)\) have under various assumptions on self-adjoint \(X_1,\ldots,X_n\)? The ambient probability space, admissible polynomials, additional assumptions, and requested regularity properties are not specified.

One conventional framework, stated conditionally rather than attributed as an explicit source hypothesis, is the following. Let \(n\geq 1\), let \((\mathcal M,\tau)\) be a von Neumann algebra with a faithful normal tracial state, and let \(X_1,\ldots,X_n\in\mathcal M\) satisfy \(X_j=X_j^*\). A noncommutative polynomial \(P\in\C\langle t_1,\ldots,t_n\rangle\) is a finite complex linear combination of words in the formal self-adjoint variables \(t_1,\ldots,t_n\). Set \(Y=P(X_1,\ldots,X_n)\). If \(Y=Y^*\), its spectral distribution is the Borel probability measure on \(\R\) defined by
\[
\mu_Y(B)=\tau\bigl(E_Y(B)\bigr)\qquad\text{for every Borel set }B\subseteq\R,
\]
where \(E_Y\) is the spectral measure of \(Y\). Equivalently,
\[
\int_{\R}f(t)\,\mathrm d\mu_Y(t)=\tau(f(Y))
\]
for every bounded Borel function \(f:\R\to\C\). Self-adjointness of the inputs alone does not ensure self-adjointness of \(Y\); the record does not impose a corresponding condition on \(P\).

In this conditional framework, the question asks for hypotheses on the tuple and the polynomial that imply specified regularity properties of \(\mu_Y\). For example, non-atomicity means \(\mu_Y(\{a\})=0\) for every \(a\in\R\), and absolute continuity means \(\mu_Y\ll\lambda\), where \(\lambda\) denotes Lebesgue measure on \(\R\). These definitions clarify possible regularity targets; they do not assert that either property holds under the assumptions displayed above.

The supplied background additionally describes a report conjecture that finite free entropy is preserved for general polynomials, and describes results about non-atomicity and other regularity under strong free-entropy-dimension hypotheses. It supplies neither an exact entropy conjecture nor the hypotheses of those results. In particular, it does not specify the entropy convention, the entropy assumption on the input tuple, the admissible output polynomials, or exclusions of constant evaluations. These omissions prevent a precise entropy implication from being stated without adding unsupported hypotheses. No freeness assumption on the inputs is explicitly supplied.
\subsection{Short English statement}
Determine which assumptions on self-adjoint noncommutative random variables and a polynomial ensure regularity of the polynomial's distribution. The supplied question leaves the assumptions and regularity targets open; interpreting the output as a probability measure on \(\R\) requires self-adjointness of the evaluated polynomial. The background also mentions an incompletely specified finite-free-entropy conjecture. The supplied label is \emph{partially solved}, a dated curation assessment requiring independent review.
\subsection{Sources}
\begin{itemize}
\item Supplied local UnsolvedMath registry record, permanent ID \(30002893\), title \emph{Regularity of Polynomial Laws in Free Probability}, exact source problem code \texttt{catalogue entry}, in the Oberwolfach Reports Open Problems collection. The record gives proposed year 2015 and names no individual proposer.
\item Supplied report citation: \emph{Free Probability Theory (2016)}, Oberwolfach Reports, pp. 1571--1629, DOI \texttt{10.4171/owr/2015/28}, \url{https://doi.org/10.4171/owr/2015/28}. This is also the supplied literature source labelled ``Oberwolfach report 28.'' The supplied citation's year 2016 and the record's proposed year 2015 are retained without resolving their relationship.
\item Supplied literature source labelled ``Charlesworth--Shlyakhtenko, regularity of free polynomials'': \url{https://www.sciencedirect.com/science/article/pii/S0022123616300830}.
\item Supplied literature source labelled ``Regularity of Polynomials in Free Variables'': \url{https://arxiv.org/abs/1408.0580}.
\item Literature assessment date: \texttt{2026-08-21}. This is supplied dated curation, not an independently verified assessment. The supplied status is ``partially solved.'' The assessment states that non-atomicity and other regularity properties are proved under strong free-entropy-dimension hypotheses, while the report conjecture concerning preservation of finite free entropy for general polynomials remains unresolved. The supplied background expressly requests independent verification of time-sensitive status and cited claims.
\item The supplied statement status is \texttt{corrected\_verified}; its accompanying verification describes only safe notation, typography, and TeX-delimiter cleanup, with hypotheses, quantifiers, and subquestions preserved. This metadata does not constitute mathematical verification of the formulation or literature claims.
\end{itemize}
\end{document}
